%% start of file `template.tex'.
%% Copyright 2006-2013 Xavier Danaux (xdanaux@gmail.com).
%
% This work may be distributed and/or modified under the
% conditions of the LaTeX Project Public License version 1.3c,
% available at http://www.latex-project.org/lppl/.


\documentclass[11pt,a4paper,roman]{moderncv}        % possible options include font size ('10pt', '11pt' and '12pt'), paper size ('a4paper', 'letterpaper', 'a5paper', 'legalpaper', 'executivepaper' and 'landscape') and font family ('sans' and 'roman')

% modern themes
\moderncvstyle{banking}                            % style options are 'casual' (default), 'classic', 'oldstyle' and 'banking'
\moderncvcolor{blue}                                % color options 'blue' (default), 'orange', 'green', 'red', 'purple', 'grey' and 'black'
%\renewcommand{\familydefault}{\sfdefault}         % to set the default font; use '\sfdefault' for the default sans serif font, '\rmdefault' for the default roman one, or any tex font name
\nopagenumbers{}                                  % uncomment to suppress automatic page numbering for CVs longer than one page

% character encoding
\usepackage[utf8]{inputenc}
\usepackage{fontawesome}
\usepackage{tabularx}
\usepackage{ragged2e}
\usepackage{tikz}
\usepackage{url}
 \usepackage{amsmath}
% if you are not using xelatex ou lualatex, replace by the encoding you are using
%\usepackage{CJKutf8}                              % if you need to use CJK to typeset your resume in Chinese, Japanese or Korean

% adjust the page margins
\usepackage[scale=0.84]{geometry}
\usepackage{multicol}
%\setlength{\hintscolumnwidth}{3cm}                % if you want to change the width of the column with the dates
%\setlength{\makecvtitlenamewidth}{10cm}           % for the 'classic' style, if you want to force the width allocated to your name and avoid line breaks. be careful though, the length is normally calculated to avoid any overlap with your personal info; use this at your own typographical risks...

\usepackage{import}

% personal data
\name{\href{https://elias-ramzi.github.io/}{Elias Ramzi}}{\color{gray} | \LARGE{ AI Research Scientist} }
% \title{Curriculum Vitae}                               % optional, remove / comment the line if not wanted
\address{Paris, France}{}{}% optional, remove / comment the line if not wanted; the "postcode city" and and


\newcommand*{\customcventry}[7][.25em]{
  \begin{tabular}{@{}l}
    {\bfseries #4}
  \end{tabular}
  \hfill% move it to the right
  \begin{tabular}{l@{}}
     {\bfseries #5}
  \end{tabular} \\
  \begin{tabular}{@{}l}
    {\itshape #3}
  \end{tabular}
  \hfill% move it to the right
  \begin{tabular}{l@{}}
     {\itshape #2}
  \end{tabular}
  \ifx&#7&%
  \else{\\%
    \begin{minipage}{\maincolumnwidth}%
      \small#7%
    \end{minipage}}\fi%
  \par\addvspace{#1}}

\newcommand*{\customcvproject}[4][.25em]{
%   \vfill\noindent
  \begin{tabular}{@{}l}
    {\bfseries #2}
  \end{tabular}
  \hfill% move it to the right
  \begin{tabular}{l@{}}
     {\itshape #3}
  \end{tabular}
  \ifx&#4&%
  \else{\\%
    \begin{minipage}{\maincolumnwidth}%
      \small#4%
    \end{minipage}}\fi%
  \par\addvspace{#1}}

% breakable skill group: renders title + content WITHOUT a minipage, so the list
% can flow across the page break and fill the page (avoids the big gap an atomic
% minipage leaves when it doesn't fit in the remaining space)
\newcommand*{\skillgroup}[2]{%
  \par\addvspace{.5em}{\bfseries #1}\par\nopagebreak
  {\small #2}\par\addvspace{.3em}}

\setlength{\tabcolsep}{12pt}

% let TeX stretch a tight line rather than push a long URL (e.g. an arxiv link) past the margin
\setlength{\emergencystretch}{3em}


\newcommand{\ExternalLink}{%
    \tikz[x=1.2ex, y=1.2ex, baseline=-0.05ex]{%
        \begin{scope}[x=1ex, y=1ex]
            \clip (-0.1,-0.1)
                --++ (-0, 1.2)
                --++ (0.6, 0)
                --++ (0, -0.6)
                --++ (0.6, 0)
                --++ (0, -1);
            \path[draw,
                line width = 0.5,
                rounded corners=0.5]
                (0,0) rectangle (1,1);
        \end{scope}
        \path[draw, line width = 0.5] (0.5, 0.5)
            -- (1, 1);
        \path[draw, line width = 0.5] (0.6, 1)
            -- (1, 1) -- (1, 0.6);
        }
    }

%----------------------------------------------------------------------------------
%            content
%----------------------------------------------------------------------------------
\begin{document}
%\begin{CJK*}{UTF8}{gbsn}                          % to typeset your resume in Chinese using CJK
%-----       resume       ---------------------------------------------------------
\makecvtitle
\vspace*{-23mm}

\begin{center}
\begin{tabular}{ c c c c }
 \faEnvelopeO\enspace \href{mailto:elias.ramzi@gmail.com}{elias.ramzi@gmail.com} & \faMobile\enspace+33\enspace6 28 25 40 88 &    \faGithub\enspace \href{https://elias-ramzi.github.io/}{elias-ramzi} & \faGraduationCap\enspace \href{https://scholar.google.com/citations?user=mbfTD_UAAAAJ}{Google Scholar} \\
\end{tabular}
\end{center}


\begin{center}
\textbf{End-to-end driving --- planning, world models, and multimodal models}
\end{center}


\section{Experiences}

{\customcvproject{valeo.ai (Paris, France) - full-time}{2024-}{
\textit{AI research scientist}\\
Deep learning research interests: end-to-end planning, world models, VLMs / LLMs.
\medbreak
\begin{itemize}
\item Developed world action models for autonomous driving, VaViM/VaVAM (technical report, CoRL Workshop, initial baseline for AlpaSim).
\item Scaled training of generative models up to 1.9B parameters across 192 H100s using $\mu$P and DeepSpeed ZeRO-2.
\item Built Pictura, a GPU-accelerated multi-agent driving simulator rendering every agent's egocentric view, which unlocks perspective-view self-play at scale for driving policies (ECCV 2026 workshop paper).
\item Co-supervising two PhD students: one on LLMs and VLMs (ICLR 2025, ACL Findings 2026), another on world models and RL.
\end{itemize}
}

{\customcvproject{Cnam \& Coexya (Paris, France) - 3 year PhD}{2021-2024}{
\textit{PhD student}\\
PhD focused on computer vision and image retrieval:
\medbreak
\begin{itemize}
    \item Designed novel ranking losses for image retrieval, combining theoretical analysis with large-scale experiments.
    \item Published four papers in major international machine learning conferences (NeurIPS, ECCV$\times$2, ICML), three of which as first author, and 2 papers in journals (TPAMI, TMLR).
    \item Transferred HAPPIER (ECCV) to trademark logo retrieval; the resulting models were deployed in production in Coexya's Acsepto product in 2024.
    \item Ran distributed training experiments across up to 16 GPUs on France's Jean Zay HPC cluster.
\end{itemize}}
}

{\customcvproject{Cnam (Paris, France) - 5 month internship}{2020}{
\textit{Research intern}\\Deep learning for 3D medical image segmentation. I used a learned confidence measure (ConfidNet) to combine segmentation results from different 2D U-Nets.
}

{\customcvproject{Sancare (Paris, France) - 6 month internship}{2019}{\textit{Data engineer and data scientist intern}
%Deep learning applied to healthcare. I participated in the development of a data pipeline to extract data from the hospitals' servers to Sancare's data model. I worked on the interpretability of a deep neural network using the layer-wise relevance propagation.
}

 {\customcvproject{Balto (Sydney, Australia) - 6 month internship}{2018}{\textit{Data engineer intern}
 %Startup working on last mile delivery of fresh food. I worked on the automation of data processing.
   % \begin{itemize}
   % \item  Implementing software :
   % \begin{itemize}
   %   \item  Python libraries to create, edit and send contractors and clients' invoices (\textasciitilde600 AUD\$ per week).
   %   \item Python libraries creating optimised delivery itineraries (\textasciitilde200 itineraries per week).
   %   \item Integration of those libraries to a graphical interface.
   % \end{itemize}
   % \item Pipeline to automate the detection of street numbers on images (this project was not finished during my internship).
 % \end{itemize}
 }}


\section{Education}
{\customcvproject{PhD in Computer Science, Cnam \& Coexya (France)}{2021-2024}
  {Ranking losses and hierarchical learning for image retrieval.
  Awarded the ``Prix de Thèse'' by \href{http://afrif.irisa.fr/?page_id=54}{AFRIF \ExternalLink} (June 2025). 6 publications including 4 first-author papers.
  \\Supervisor: \href{https://thome.isir.upmc.fr}{Nicolas Thome \ExternalLink} (Sorbonne Université).
  \\Advisors: \href{https://nicolas.audebert.at/}{Nicolas Audebert \ExternalLink} (IGN); \href{https://clementrambour.github.io}{Clément Rambour \ExternalLink} (Sorbonne Université).
  \\Industrial advisor: \href{https://orcid.org/0009-0006-7254-6178}{Xavier Bitot \ExternalLink} (Coexya).
  }
}
{\customcvproject{Engineering diploma, CentraleSupélec (France)}{2016-2020}
  {Courses in signal processing (sound, image, speech), machine learning, statistical models (estimators, Bayesian learning), data science (data mining, sparse data representation, NLP).}
}

%{\customcvproject{Preparatory school, PSI*, Lycée Lakanal (France)}{2014-2016}
%  {Intense courses in mathematics and physics.}
%}


\section{Technical skills}
\skillgroup{Open source projects}{
   \begin{itemize}
   \setlength{\itemsep}{0pt}
   \item Pictura simulator: \href{https://github.com/valeoai/Pictura/}{\faGithub\enspace Pictura}
   \item Overleaf MCP server: \href{https://github.com/elias-ramzi/WebLatexMCP/}{\faGithub\enspace WebLatexMCP}
   \item Reproduce VaViM and VaVAM (CoRL'25): \href{https://github.com/valeoai/VideoActionModel/}{\faGithub\enspace VideoActionModel}
   \item Global and local prompts for VLMs (ECCV'24): \href{https://github.com/MarcLafon/gallop/}{\faGithub\enspace GalLoP}
   \item Out-of-distribution detection with HEAT (ICML'23): \href{https://github.com/MarcLafon/heatood/}{\faGithub\enspace HEAT}.
   \item Ranking losses and hierarchical image retrieval (NeurIPS'21, ECCV'22, TPAMI'25): \href{https://github.com/elias-ramzi/ROADMAP}{\faGithub\enspace ROADMAP}, \href{https://github.com/elias-ramzi/HAPPIER}{\faGithub\enspace HAPPIER}, \href{https://github.com/elias-ramzi/SupRank}{\faGithub\enspace SupRank} \& the hierarchical \href{https://github.com/cvdfoundation/google-landmark}{\faGithub\enspace google-landmark} dataset.
 \end{itemize}}

\skillgroup{Technologies}{
   \begin{itemize}
   \setlength{\itemsep}{0pt}
   \item  {My programming language:} Python -- PyTorch, Lightning, NumPy, Hydra, Pandas, Scikit-learn.
   \item  {My work environment :} Linux/macOS + Git + VS Code \& claude-code.
   \item Other tools: \LaTeX, Shell script, SLURM, HyperQueue, Nextflow.
 \end{itemize}}

\skillgroup{Contributions to open-source}{
   \begin{itemize}
   \item PyTorch metric learning: \href{https://github.com/KevinMusgrave/pytorch-metric-learning/pull/303}{\faGithub\enspace pytorch-metric-learning}.
   \item Implementation of Mixture of Experts: \href{https://github.com/davidmrau/mixture-of-experts/pull/14}{\faGithub\enspace mixture-of-experts}.
   \item PyTorch-Lightning: \href{https://github.com/Lightning-AI/lightning/pull/2319}{\faGithub \enspace lightning}.
   \item Universal Image Embeddings: \href{https://github.com/nikosips/Universal-Image-Embeddings/tree/master\#acknowledgements}{\faGithub \enspace  Universal-Image-Embeddings}.
 \end{itemize}}


\section{Miscellaneous}

\begin{itemize}
    \item French is my mother tongue and I speak fluently English (last tested at C1 level).
    \item I enjoy biking (\href{https://github.com/elias-ramzi/ClaudeCyclingMCP/}{\faGithub\enspace ClaudeCyclingMCP}), (trail) running and back-country skiing.
\end{itemize}


 \section{Publications}

Yuan Yin, \textbf{Elias Ramzi}, Marc Lafon, Valentin Charraut, Victor Bares, Yihong Xu, {\'E}loi Zablocki, Alexandre Boulch, Thibault Buhet, Andrei Bursuc, \textit{et al.} \href{https://arxiv.org/abs/2607.26005}{"Pictura: Perspective-View Self-Play at Scale for Driving."} \textit{ECCVw}, 2026.

 \smallskip

Yihong Xu, {\'E}loi Zablocki, Yuan Yin, \textbf{Elias Ramzi}, Ellington Kirby, Alexandre Boulch, Matthieu Cord. \href{https://arxiv.org/abs/2606.07170}{"TOAD: Test-Time Trajectory Optimization for Autonomous Driving."} \textit{arXiv preprint}, 2026.

 \smallskip

Shashanka Venkataramanan, Valentinos Pariza, Mohammadreza Salehi, Lukas Knobel, Spyros Gidaris, \textbf{Elias Ramzi}, Andrei Bursuc, Yuki M. Asano, \href{https://arxiv.org/abs/2507.14137}{"Franca: Nested Matryoshka Clustering for Scalable Visual Representation Learning."} \textit{CVPR}, 2026.

 \smallskip

 Amaia Cardiel, {\'E}loi Zablocki, \textbf{Elias Ramzi}, and Eric Gaussier. \href{https://arxiv.org/abs/2603.00696}{"DRIV-EX: Counterfactual Explanations for Driving LLMs."} \textit{ACL findings}, 2026.

  \smallskip

 Florent Bartoccioni, \textbf{Elias Ramzi}, \textit{et al.} \href{https://arxiv.org/abs/2502.15672}{"VaViM and VaVAM: Autonomous Driving through Video Generative Modeling."} \textit{Technical report \& CoRL workshop}, 2025.

\smallskip

Amaia Cardiel, {\'E}loi Zablocki, \textbf{Elias Ramzi}, Oriane Sim{\'e}oni, and Matthieu Cord. \href{https://arxiv.org/abs/2409.11919}{"LLM-wrapper: Black-Box Semantic Aware Adaptation of Vision-Language Models for Referring Expression Comprehension."} \textit{ICLR}, 2025.

  \smallskip

     \textbf{Elias Ramzi}, Nicolas Audebert, Cl{\'e}ment Rambour, Andr{\'e} Araujo, Xavier Bitot and Nicolas Thome. \href{https://arxiv.org/abs/2309.08250}{"Optimization of Rank Losses for Image Retrieval."} \textit{TPAMI}, 2025.

\smallskip

 Marc Lafon, \textbf{Elias Ramzi}, Cl{\'e}ment Rambour, Nicolas Audebert, Nicolas Thome. \href{https://arxiv.org/abs/2407.01400}{"GalLoP: Learning Global and Local Prompts for Vision-Language Models."} \textit{ECCV}, 2024.

 \smallskip

Yannis Karmim, \textbf{Elias Ramzi}, Raphaël Fournier S’niehotta, and Nicolas Thome. \href{https://arxiv.org/abs/2407.07912}{"ITEM: Improving Training and Evaluation of Message-Passing based GNNs for top-k recommendation."} \textit{TMLR}, 2024.

 \smallskip

  Marc Lafon, \textbf{Elias Ramzi}, Cl{\'e}ment Rambour, Nicolas Thome. \href{https://arxiv.org/abs/2305.16966}{"Hybrid Energy Based Model in the Feature Space for Out-of-Distribution Detection."} \textit{ICML}, 2023.

  \smallskip

 \textbf{Elias Ramzi}, Nicolas Audebert, Nicolas Thome, Cl{\'e}ment Rambour, and Xavier Bitot. \href{https://arxiv.org/abs/2207.04873}{"Hierarchical Average Precision Training for Pertinent Image Retrieval."} \textit{ECCV}, 2022.

\smallskip

 \textbf{Elias Ramzi}, Nicolas Thome, Cl{\'e}ment Rambour, Nicolas Audebert, and Xavier Bitot. \href{https://arxiv.org/abs/2110.01445}{"Robust and Decomposable Average Precision for Image Retrieval."} \textit{NeurIPS}, 2021.



\end{document}
